%!TEX root = ../template.tex
%%%%%%%%%%%%%%%%%%%%%%%%%%%%%%%%%%%%%%%%%%%%%%%%%%%%%%%%%%%%%%%%%%%
%% chapter4.tex
%% NOVA thesis document file
%%
%% Chapter with introduction
%%%%%%%%%%%%%%%%%%%%%%%%%%%%%%%%%%%%%%%%%%%%%%%%%%%%%%%%%%%%%%%%%%%

\typeout{NT FILE chapter4.tex}%

\chapter{Evaluation}
\label{cha:evaluation}

% epigraph configuration
\epigraphfontsize{\small\itshape}
\setlength\epigraphwidth{12.5cm}
\setlength\epigraphrule{0pt}
%
%
%
\prependtographicspath{{Chapters/Figures/chapter4}}

In this chapter, we evaluate the developed solution, while displaying how it performs with a case study. Firstly, we begin by explaining our case study, how it was implemented, providing the "point of view" of a developer using the developed Overlord framework. In Section~\ref{ssec:monitoring_reliability}, we detail the role that average reliability plays in providing a basis for assessing the system's overall performance. Section~\ref{ssec:fanout_adaptation_process} describes the process of adapting the fanout of the broadcast protocol according to the monitored reliability through the use of developer designed rules. Finally, In Section~\ref{ssec:evaluation_scenarios}, we evaluate the developed solution under different scenarios that might occur during the execution of a system like this, testing the adaptive application against static fanout alternatives.

\section{Message Sharing Application}
\label{sec:message_sharing_application}
The case study used for evaluation of our solution is a message sharing application. Each participant of the system generates messages periodically, and the main objective of the application is for these messages to reach all other participants. For this purpose, the application uses a gossip-based broadcast protocol in order to disseminate messages through the system. As discussed in Section~\ref{ssec:gossip_protocols}, in a gossip-based broadcast protocol, each node, when it receives a message, forwards it to $x$ of its neighbors, and depending on the value of $x$, the message will be received by every participant of the system with a high probability. This value is denominated as \textit{fanout}, and as also discussed in Section~\ref{ssec:membership}, when the fanout is equal to $\log{n}$, where $n$ is the number of participants in the system, the message has a high probability of reaching every node (assuming there are no kinds of disturbances occurring)\cite{Leitao2010a, Kermarrec2003}.

Since we are dealing with a pure decentralized system, it should be dynamic, meaning that the number of participants should be able to change. To ensure that all participants in the system receive all the messages, the fanout parameter has to be adjusted in relation to the number of nodes. Moreover, if some messages are lost, or some of the nodes do not forward a message, increasing the fanout should also increase the probability of all nodes receiving every message. Finally, if the number of participants reduces drastically and the fanout is too high for this new number of participants, the number of duplicate messages will also increase. Managing the fanout parameter will be the main task of Overlord in order to ensure that, even with disturbances in the system, a high percentage of the participants will receive all the messages.

\paragraph{The Role of Reliability} \label{par:the_role_of_reliability} As introduced in \ref{par:reliability}, the reliability of a system is an important aspect concerning the performance of our system. Consequentially, it directly determines the fanout value decided by Overlord. This means that monitoring the reliability of the system is a crucial component, as it determines if the main objective is being fulfilled. In the context of our experiments, it was decided that a reliability above 98\%\todo{Check if this is still true for all experiments.} is acceptable, although this value can change and is a configurable parameter, that will be referred to as $l$ in the remainder of this section, in order to keep the specifications of this use case generic and independent of any concrete values.

\subsection{Monitoring Reliability}
\label{ssec:monitoring_reliability}
As discussed before, Overlord tasks each node with monitoring their protocols locally, and periodically, and makes use of MON-Collect to transport and aggregate this monitored data to the Overlord central node. In the context of this case study, this means that each node must monitor the broadcast protocol. This protocol implements an adaptive version of an existing eager push gossip broadcast protocol\cite{Leitao2010a}. This implementation is adaptive, which means that, the fanout parameter can be changed upon request, which is sent by Overlord when decisions are made by the Rule Engine. Moreover, this implementation also had to be expanded in order to include metric monitoring. For this purpose, we used the metrics framework available in Babel\cite{Paulico2025}.

This version of the broadcast protocol enables each node to monitor its individual performance, which can be aggregated and processed with others' monitored information at the Overlord node. The broadcast protocol monitors every message, creating a record for each of these messages. For each recorded message, the broadcast protocol stores different information:
\begin{itemize}
    \item \textbf{Sent Messages:} A sent message record contains the ID of the message (its unique identifier), the time in milliseconds at which it was created, and its creator's identifier.
    \item \textbf{Received and Delivered Messages:} These records contain the message ID, the identifier of its creator, the timestamp of its creation, the time at which it was delivered, and its hop count (as explained in Section~\ref{ssec:gossip_protocols}).
\end{itemize}

\paragraph{In-Node Aggregation} \label{par:in_node_aggregation} When MON-Collect triggers a collection, each node will first request the monitored metric information from all the appropriate protocols. In the context of this case study, the only protocol is the broadcast protocol, which will then provide all the records it created about the messages sent and received since the last metric collection. Using the Babel metrics framework, it is possible to perform what we refer to as an in-node aggregation. Essentially it serves as a way for processing the raw monitored information from the broadcast protocol into usable and meaningful information. The in-node aggregation is responsible for creating \textbf{message records}, which are used to facilitate in-network aggregation of message information between nodes.

\paragraph{Message Records} \label{par:message_record} In-node aggregation processes different records of the same a message into a single record. This record contains all the necessary information about a message that Overlord needs in order to make decisions. A message record contains: \begin{inparaenum}[(i)] \item the message's unique identifier, \item its creation timestamp, \item the highest recorded delivery timestamp, \item the number of times the message was delivered, \item the number of times it was received, \item the number of duplicates and \item its LDH (last delivery hop).\end{inparaenum}

Raw records created from the broadcast protocol are transformed into message records. When MON-Collect collects the monitored information from all other nodes, it joins the message records into the same single record. This way, only the important information is stored and transmitted to the next nodes, as unnecessary information is discarded.

Figure~\ref{fig:recordAggregation} illustrates this phenomenon. Figure~\ref{fig:recordAggregation1} presents an alternative where no aggregations are made. As both node $b$ and $c$ receive the message with ID 1, both create a message record. While node $b$ delivered the message at $t1$, node $c$ delivered the message at time $t2$, where $t1 > t2$. Moreover, in the case of node $b$, the message hop count is $h1$, while for node $c$, the message hop count is $h2$, where $h2 > h1$. Both nodes send this record to node $a$, their parent. However, without any aggregations, node $a$ simply forwards this data to its parent. This is not desirable since, in this case, Overlord only needs the message's highest delivery timestamp, and its LDH. Figure~\ref{fig:recordAggregation2} presents the alternative, where instead of simply forwarding the two records, it aggregates them into a new $m3$ message record, discarding both $t2$ and $h1$, keeping only the highest values for that message.

\begin{figure}
\centering
\begin{subfigure}{.46\textwidth}
  \centering
  \includegraphics[scale=1.25]{recordAggregation1.pdf}
  \caption{Node $a$ forwards both $m1$ and $m2$ to its parent, keeping both hop counts and delivered timestamps.}
  \label{fig:recordAggregation1}
\end{subfigure}
\hspace{1cm}
\begin{subfigure}{.46\textwidth}
  \centering
  \includegraphics[scale=1.17]{recordAggregation2.pdf}
  \caption{Instead of forwarding both records to its parent, node $a$ creates $m3$ with only the necessary information.}
  \label{fig:recordAggregation2}
\end{subfigure}
\caption{A Visual representation of the benefits of record aggregation in this case study}
\label{fig:recordAggregation}
\end{figure}

With this information about each message, the Overlord is able to calculate the necessary metrics to make any decisions about adapting the fanout, as well as some other relevant metrics that will help a developer understand how the system is performing (for instance, the reliability might be high, however message latency can still be high, and it can be important for a developer to know this\todo{I don't think I'm explaining myself very well here.}). When a collection is taking place, and a node has received its own monitored data, as well as other nodes', this node will aggregate each message record, keeping only the LDH, the highest timestamp of delivery, and summing all the delivery, received and duplicate counts. As the nature of MON-Collect is to create a tree for collection, it ensures that no loops are created without any extra steps, meaning that the information of one single node is considered exactly once, making it possible for these summations and aggregations to be done without any risks.

\subsection{Fanout Adaptation Process}
\label{ssec:fanout_adaptation_process}

With the metric collection of every node of the system at the Overlord node, it must detect if the system is showing good performance, and if not, how to counter these issues by means of adapting the fanout parameter. For this purpose we use the Rule Engine, the behavior of which can be found more detailed in Section~\ref{sec:ruleengine}.

Since Overlord makes decisions through \gls{eca} rules, one must be created to deal with values of reliability below $l$ (the threshold of acceptable reliability). The first iteration of this rule can be found at Pseudocode~\ref{code:reliabilityv1}.
\begin{figure}[ht]
    \centering
    \begin{algorithm}[H]
        \caption{Handle Reliability Rule - First Iteration}
        \label{code:reliabilityv1}
        \begin{algorithmic}
            \IF{$reliability < l$}
                \STATE $fanout \gets fanout + 1$
            \ENDIF
        \end{algorithmic}
    \end{algorithm}
\end{figure}

Although the reliability directly affects the decided fanout, it is not the only consideration to be made. Having a fanout unnecessarily high introduces a lot of duplicate messages, which is undesirable if they are not necessary. If increasing the fanout does not improve the reliability, then we shouldn't increase it infinitely, as it risks polluting the network. Because of this, we introduced a cap value to the fanout, which is defined as $\log{n} + c$, where $n$ is the number of nodes in the system, and $c$ is a configurable parameter. Like this, the fanout can increase to help with any system disturbance, however it cannot grow infinitely. This cap can hinder possible improvements from Overlord, since it limits how much it can increase the fanout. There might be a situation where increasing the fanout more than $c$ might produce a better reliability. However, this comes down to a trade-off, as it is far more dangerous to let Overlord increase this value infinitely. The developer must adjust this $c$ in a way that it allows Overlord to improve system performance, but not at the cost of a scenario where redundancy becomes too high.\todo{I added this part. Review.}

Another aspect to consider, is updating the fanout as the number of nodes changes. This is not a straightforward matter. For example, if the Overlord receives monitored information from 200 nodes, and in the next collection it only receives information from 30 nodes, it is possible that the number of participants $n$ decreased. However, it is also possible that this is the result of a failure. Moreover, even if the number of participants decreases, we shouldn't directly reduce the fanout to $\log{n}$, as it could be a drastic change to the system, and possibly make it function in a way that is undesired. Instead, we should opt for a more careful transition.

Moreover, decreasing the fanout when the reliability is only acceptable, might not be desirable, as this may produce an even lower reliability, possibly under the threshold $l$. Consequentially, the fanout will only decrease if the reliability is at 100\%. As such, when the node count of the system decreases nothing happens. Only when $d$ consecutive collections are made (where $d$ is a configurable parameter),  where $\log{n}$ remains unchanged and the reliability is always at 100\%. Decreasing the fanout happens one by one until the fanout is equal to $\log{n}$, or until the reliability goes bellow the $l$ defined threshold.

Combining all aspects explained, and considering that both $c$, $d$ and $l$\todo{Maybe change the letters out for actual variable names} are configurable parameters, the pseudocode that defines the reliability rule is presented in Pseudocode~\ref{code:reliabilityv2}.
\begin{algorithm}[H]
    \caption{Handle Reliability Rule}
    \label{code:reliabilityv2}
    \begin{algorithmic}[H]
        \IF{$reliability < l$}
            \STATE $counter \gets 0$
            \IF{$fanout + 1 < \log{n} + c$}
                \STATE $fanout \gets fanout + 1;$
            \ENDIF
        \ELSIF{$\log{last\_n} < \log{n}$}
            \STATE $counter \gets 0$
            \STATE $fanout \gets fanout + 1;$
        \ELSIF{$\log{last\_n} > \log{n}$ and $reliability = 1$}
            \IF{$counter = d$}
                \STATE $counter \gets 0;$
                \STATE $fanout \gets fanout - 1;$
            \ELSE
                \STATE $counter \gets counter + 1;$
            \ENDIF
        \ENDIF
    \end{algorithmic}
\end{algorithm}

\subsection{Evaluation Scenarios}
\label{ssec:evaluation_scenarios}

With a rule to adapt the fanout parameter in place, we require scenarios to test how it performs and determine its benefits in relation to a solution without Overlord. For this purpose, we designed a few scenarios that allow us to evaluate Overlord's performance against a solution without Overlord. These scenarios are:
\begin{inparaenum}[(i)]
    \item Node influx, in which we test how Overlord adapts the fanout upon new nodes entering the system,
    \item Crashed nodes or leaving nodes, which is a scenario that evaluates how Overlord acts when the amount of participants of the system reduces,
    \item Free riders scenario, which tests if increasing the fanout parameter to a value higher than $\log{n}$, improves reliability in case some nodes choose to only receive messages and not forward them to neighbors, and
    \item Packet loss, which allows us to test this same aspect, only instead of some nodes not forwarding messages, a certain percentage of messages are lost and not delivered.
\end{inparaenum}

\paragraph{Overlord Failure} Another aspect we considered while testing and evaluating Overlord, was the possible failure of a node during the creation of the tree. Since Overlord collects metrics through a tree, if one node fails and does not respond, then all its collected data will not be received by Overlord. Moreover, any children of this node will not have its collected data received by Overlord. The closer to Overlord the node that does not respond is, the more nodes will be affected. To evaluate this aspect, we will also perform an evaluation with different percentages of failed nodes and determine Overlord's behavior under the free riders scenario.

\subsubsection{Node Influx}
\label{sssec:node_influx_experiment}
In this scenario, the system will experience an influx of new participants. The objective is to observe how Overlord adapts the fanout parameter in relation to the number of participants joining the system, and how this compares to a system with a static fanout.

The same experiment will be conducted for two\todo{Change this to 3 when ready.} versions of the application. One with Overlord, and another without it. Groups of new participants will be joining periodically. The number of new participants in each period corresponds to the number needed for Overlord to increase the fanout parameter. Table~\ref{table:nodes_vs_fanout}, presents the relation between the number of nodes in the system, and their expected optimal fanout\cite{Leitao2007a}.

\begin{table}
    \centering
    \begin{tabular}{|c|c|}
    \hline
    Number of Nodes & Expected Fanout \\ \hline
    20              & 3               \\ \hline
    54              & 4               \\ \hline
    148             & 5               \\ \hline
    403             & 6               \\ \hline
    \end{tabular}
    \caption{Number of nodes in the system vs the expected fanout the broadcast should be set to.}
    \label{table:nodes_vs_fanout}
\end{table}

\paragraph{Experiment Details} \label{par:node_influx_experiment_details} The experiment spanned 1 hour and 40 minutes, having 4 district periods, in which new participants join. It begins with 20 nodes. Around minute 17, 34 new nodes join. The second period of influx begins at minute 35, where 94 new nodes join, which means a total of 148 nodes in the system. Just after the 57-minute mark, the final wave of nodes starts joining. 262 nodes start joining the system, making our final amount of nodes to be 410 (just over the 403 needed). As the amount of new nodes here is significant, this process takes around 20 minutes. As such, after all nodes join, the system runs without interference for an additional 25 minutes.

\begin{figure}
    \centering
	\includegraphics[scale=0.93]{NodeInflux/nodesEvolution.pdf}
	\caption{Node count evolution over time}
	\label{fig:node_influx_nodes_evolution}
\end{figure}

Figure~\ref{fig:node_influx_nodes_evolution} presents the evolution of the number of nodes in the system, as well as the decisions made by Overlord. In order to understand such decisions, it is also necessary to understand the monitored reliability of the system during periods where large amounts of nodes join.

Figure~\ref{fig:node_influx_overlord_reliability_evolution} presents this information. As it is observable, the reliability lowers during periods where new participants are joining the system. This is a consequence of older messages still being processed when new nodes join. For instance, consider a message that is created and disseminated. This message reaches all participants of the system. Before the next collection of metrics, new participants join, and start receiving new messages. However, during the next collection of the system, Overlord has no way of knowing how many nodes were in the system for each message (as this would imply a more complex approach to monitoring metrics). Even with a solution, there is no clear way to decide if a newer node should receive messages sent before it joined the system\todo{I really hate this because I just thought of a simpler solution, but oh well...}. Consequentially, when processing messages disseminated previously to new nodes joining the system, the Overlord assumes these nodes should have received the message. This presents us with the lower reliability observed.

Consequentially, we observe that Overlord tries to improve the reliability by increasing the fanout parameter. For instance, during the second influx of nodes, the Overlord increases the fanout to 6 and 7, as a consequence of lower reliability, not of the number of nodes increasing. However, just before nodes start joining once more, Overlord starts the optimization phase, checking that reliability is acceptable, and decreasing the fanout to 6. If the new wave of nodes had not joined, Overlord would have decreased the fanout once more, reaching its optimal state. 

Moreover, between minute 60 and minute 80, the reliability is constantly low, as the Overlord reaches the maximum fanout it can increment to. Overlord tries to increase the fanout due to the low reliability, however it does not go above 9, since it reaches its maximum incremental value during this wave of new nodes. This restriction would only be removed when the system reached the 403 number of participants, where Overlord could increase the fanout once more to 10, as $\log{n}$ would have also increased. However, as reliability increases at the end of this wave, Overlord begins the optimization phase, reducing the fanout. Had the experiment ran for longer, the fanout would have reached its optimal state.

\begin{figure}
    \centering
	\includegraphics[scale=0.93]{NodeInflux/overlordReliabilityEvolution.pdf}
	\caption{Average reliability of Overlord evolution over time.}
	\label{fig:node_influx_overlord_reliability_evolution}
\end{figure}

This phenomenon is undesired, and is mostly a limitation of choosing to monitor metrics on demand and in a centralized manner\todo{Since I now have come up with a solution, I feel like I'm just lying here}. However, considering the bigger picture, the rule we designed should include all possible scenarios that could occur during the application's lifespan. This means that the rule has no way of knowing if the reliability is low as a consequence of the increase of nodes, or as a consequence of another occurring phenomenon, so it decides to take a cautious approach, and increase fanout nonetheless to help reliability. The cap to the increase of the fanout parameter was created as to, in the event that Overlord makes the wrong decision, and increasing the fanout does not improve reliability, the cost of this mistake does not grow infinitely, and is only in place during this disturbance in the reliability. When the system reaches higher reliability, then it begins the phase of optimization, as the system is no longer underperforming. As a consequence of improving system performance, it does not always make the correct decision, however, it always tends to the most optimal state\todo{I really do not like this part. Should I just say that I didn't have time to implement the correct solution, but the behavior makes sense and improves the system?}.

\begin{figure}
    \centering
	\includegraphics[scale=0.93]{NodeInflux/reliabilityComparisonEvolution.pdf}
	\caption{Comparison of the average reliability evolution over time.}
	\label{fig:node_influx_reliability_comparison_evolution}
\end{figure}

After discussing these results, it is important to compare them to a version of this application not running Overlord, with a static configuration, which is depicted in Figure~\ref{fig:node_influx_reliability_comparison_evolution}. Overlord presents the same reliability curve in Figure~\ref{fig:node_influx_overlord_reliability_evolution}, while adding the curve correspondent to application with a static optimistic fanout of 3 (considering the initial state where only 20 nodes are present). Both curves are similar, where Overlord experiences marginal improvements in relation to the static application after the increase in number of participants. The Overlord application, when no new participants are joining the system, reports a reliability of around 100\%, where the optimistic application reports a reliability of 99\%. Although this result does not portray a significant improvement, a reliability of 99\% indicates that, for every message, at least one node does not receive it, which contradicts the main objective of the application.

Moreover, upon normal conditions of the system, Overlord will always tend to adapt the fanout to be $\log{n}$, which will, with a high probability, mean that every node of the system will receive each message\cite{Leitao2007a}. Even though under normal conditions, the optimistic application does not present significant underperformance, the experiments in Sections~\ref{sssec:free_riders} and \ref{sssec:packet_loss}, will show how, even with the most optimal fanout value (relatively to the number of participants), when the system is under disturbances, its performance is significantly affected. Consequentially, having a fanout value even lower than the optimal one, will only increase this dip in performance.

Thus, our solution tends to correctly adapt the fanout parameter to each influx of nodes, while also proving itself more resilient, as the optimistic version of the application, not only underperforms under normal conditions, while also being extremely vulnerable to any disturbances that might occur in the system.

\subsubsection{Crashed Nodes}
\label{sssec:crashed_nodes_experiment}
In this experiment, we test how Overlord reacts to the scenario where some nodes of the system crash. Similarly to the experiment in Section~\ref{sssec:node_influx_experiment}, we observe how Overlord adapts the fanout parameter to the shifting number of participants in the system, however, in this case, instead of this number increasing, it decreases. It also has the component of crashing the nodes, meaning that they leave the system immediately on command. The main objective with this scenario, is to observe how Overlord optimizes the fanout parameter, decreasing it when the system is performing correctly.

\paragraph{Experiment Details} \label{par:crashed_nodes_experiment_details} This experiment ran for a total of 83 minutes. The amount of nodes crashed evolves through waves, and in each wave, the number of nodes that leaves the system, is the exact number necessary to reduce the fanout (as dictated by the $\log{n}$ calculation), as depicted in Table~\ref{table:nodes_vs_fanout}, meaning that with each wave, we should observe Overlord reduce the fanout. The experiment begins with 350 nodes, while decreasing to 148 at the 15-minute mark. At minute 30, the number of nodes drops once again to 54. The final wave of crashed nodes begins just before the 60-minute mark, where the number of nodes decreases to 20 nodes. The application then runs with the 20 nodes for 25 minutes.

\begin{figure}
    \centering
	\includegraphics[scale=0.93]{CrashedNodes/nodesEvolution.pdf}
	\caption{Node count evolution over time.}
	\label{fig:crashed_nodes_overlord_evolution}
\end{figure}

Figure~\ref{fig:crashed_nodes_overlord_evolution} illustrates the evolution of the number of nodes of the system, as well as the decisions that the Overlord made overtime. As we can observe, with every decrease in number of nodes, the Overlord, after three collections where reliability is stable, decides to optimize the fanout, and reduce it, which is the intended behavior.

An important matter, is concerning the third collection, where the reported number of nodes is just above 50, increasing then to 148. This represents a failure in the collection phase, as when the metric collection tree was being created, some nodes crashed. Both their monitored data, and the one from their children was not considered in that collection. However, in the next metric collection, the amount of nodes reported is the correct one, and it does not prevent Overlord from making the correct decision. More on how loss of information affects Overlord can be found under Section~\ref{sssec:tree_drop}.

Now we must compare the performance of our solution to static configuration alternatives. In this case, we will compare it against one that sets its fanout to 6 since it is the most optimal fanout for the number of nodes at the beginning of execution, which is the pessimistic application, as it considers only the worst scenario (in the experiment). We will also run an optimistic version with its fanout set to 3, which corresponds to the most optimal fanout for when the system has only 20 participants, which happens at the end of execution.

\begin{figure}
    \centering
	\includegraphics[scale=0.93]{CrashedNodes/reliabilityEvolutionFanout.pdf}
	\caption{Average reliability evolution over time.}
	\label{fig:crashed_nodes_reliability_evolution}
\end{figure}

Figure~\ref{fig:crashed_nodes_reliability_evolution} presents the average reliability reported by the three versions of the application tested. Both the pessimistic and Overlord applications were constant in terms of reported average reliability, reporting always around 100\%. However, the optimistic application, even though it reported high reliability during periods of normal execution, during the crash of nodes, reliability was affected, dropping to 90\% in the worst case, whereas the other alternatives did not suffer from this issue. Moreover, since we wish to compare optimization in terms of efficiency, we will be comparing their reported average RMR. 

\begin{figure}
    \centering
	\includegraphics[scale=0.93]{CrashedNodes/rmrEvolutionV2.pdf}
	\caption{Average RMR evolution over time.}
	\label{fig:crashed_nodes_rmr_evolution}
\end{figure}

This comparison can be observed in Figure~\ref{fig:crashed_nodes_rmr_evolution}. At the beginning of execution, both the pessimistic and Overlord applications report the same RMR values. Just when Overlord first decides to optimize, it immediately reports a lower RMR, which considering it keeps the same average reliability, means that it achieves the same performance, with less redundancy. This decrease in RMR happens with every optimization decision of Overlord. Consequentially, with each decision, Overlord's reported RMR begins to draw itself apart from the one reported by the pessimistic application. Towards the end of the experiment, Overlord reports an average RMR of 2, where the pessimistic application still indicates an average RMR of around 5, meaning that for every node, each message is received 3 times more than in the Overlord application, which considering their performance is the comparable, means this value is inadequate.

Moreover, the optimistic application maintains a lower RMR value, being always around 1. This value is optimal as it means that each node receives a message exactly once, with almost no duplicates being received. However, as mentioned previously, this usually comes at the cost of a high reliability value. With Figure~\ref{fig:crashed_nodes_reliability_evolution}, we can observe this trade-off, as the optimistic application experiences drops in reliability during instability phases of the system. This means that, although this application is efficient and almost generates no duplicate messages, it fails to be reliable, which is the main objective of the application. Moreover, in the next experiments, we can also observe how reliability is affected by other disturbances that can occur in the system, which will reinforce how fragile a lower fanout option is.

\subsubsection{Free Riders}
\label{sssec:free_riders}
This experiment reflects a specific scenario, in which some nodes only want the benefits of participating in the system, without contributing to their responsibilities with the intent of saving resources. Many studies reported high percentages of free-riders in \gls{p2p} systems\cite{Singh2018}\todo{This paper is very sus. It says this but does not cite anything.}. In this case, this implies that these nodes do not forward messages to their neighbors after delivering them. It impacts reliability directly, which provokes Overlord to react and try to increase it. In this experiment, free riders are referred to as "pit nodes", since this better reflects their impact on the system, rather than their intention. A certain percentage of the nodes of the system are chosen to be "pit nodes". In order to increase the reliability under these conditions, Overlord will increase the fanout of the broadcast protocol. Figure~\ref{fig:freeridersexperiment1} represents a small portion of a hypothetical system where the fanout parameter is set to 4, and nodes colored red are pit nodes, while green nodes are normal functioning nodes. In this figure, node $a$ will not receive any messages that are forwarded by any of its neighbors because all of them are pit nodes. If we increase the fanout to 5, another node that might contain the same messages can be added to its membership, which means that this new node will forward the message to node $a$ (as is shown in Figure~\ref{fig:freeridersexperiment2}). It is possible that this new node is also a pit node, although the fanout can increase again until most nodes have at least one normal node as a neighbor. Another caveat of this experiment, is concerning messages created by pit nodes, which in this case are forwarded to its neighbors.

\begin{figure}
\centering
\begin{subfigure}{.46\textwidth}
  \centering
  \includegraphics[scale=2.1]{FreeRiders/freeridersA.pdf}
  \caption{Node $a$ will not receive message $m$ since all the neighbors with message $m$ are pit nodes.}
  \label{fig:freeridersexperiment1}
\end{subfigure}%
\hspace{1cm}
\begin{subfigure}{.46\textwidth}
  \centering
  \includegraphics[scale=2.1]{FreeRiders/freeridersB.pdf}
  \caption{As the fanout increased to 5, $a$ now has a new neighbor $f$ which delivers $m$ to it.}
  \label{fig:freeridersexperiment2}
\end{subfigure}
\caption{A visual example of the expected improvements of Overlord on the free riders experiment}
\label{fig:freeridersexperiment3}
\end{figure}

\paragraph{Experiment Details} \label{par:freeridersexperiment_details} This experiment was done with 200 participants in the system. The  defined acceptable reliability is above 98\%. We will be comparing three versions of the message dissemination application: \begin{inparaenum}[(i)] \item A version with an optimistic static fanout, which does not consider any disturbances on the system, \item A version with a pessimistic static fanout, considering the consequences of this experiment, and \item A version running Overlord and the rule specified in Section~\ref{ssec:fanout_adaptation_process}, which will adapt the fanout in accordance to the monitored reliability. \end{inparaenum}For 5 metric collections the system works as indented (to simulate normal behavior). At the 5th collection, pit nodes are turned on, and about 50\% of the nodes stop forwarding messages. Pit nodes are turned on for 8 more rounds, giving time for Overlord to react and its changes to take effect and become visible.

Another aspect to consider, is the bias that randomization introduces. Since we use HyParView\cite{Leitao2007}, for different experiments, the overlay network connecting the nodes will be different. Moreover, the nodes that become pit nodes are also selected randomly. This is undesirable for comparing different solutions, because these variables can introduce a bias not related to the performance of the different solutions. Since we are comparing three different versions of the same application (optimistic static fanout, pessimistic static fanout and overlord adaptive fanout), it is important that for these different versions, the overlay networks created are the same, and the decided pit nodes are also the same, in order to reduce possible biases. This does not mean that only one membership and pit node set is used in the experiments, but for all experiments, the same sets of memberships and pit nodes were used for all versions of the application.

\begin{figure}
    \centering
	\includegraphics[scale=0.9]{FreeRiders/freeRidersOverlordReliability.pdf}
	\caption{Average reliability evolution over time}
	\label{fig:freeridersOverlordReliability1}
\end{figure}

\par In Figure~\ref{fig:freeridersOverlordReliability1}, we present the evolution of the average reliability over time of the application that is running Overlord, in combination with the different fanout values it decided. The area that is shaded with red, represents the time during which pit nodes were activated, and the vertical markers show the decided fanout by Overlord as a consequence of the metrics collected. The initial fanout is 5. From the start of the experiment to the 4-minute mark, the system runs normally, and the reliability is constantly at the 100\% mark. At the 4th minute, pit nodes are activated. The effect is already visible in the next metric collection, as reliability drops to 97\%, which is bellow the acceptable threshold. The effect becomes more visible in the next collection, as some messages sent before pit nodes activation are not yet processed at this time. At minute 6, the reliability is at its lowest, being below 94\%, which is undesirable. Although it is a high number, for a system tasked with broadcasting messages to all participants, it means that in this case, in average, for every message, 12 nodes did not receive it. This triggers Overlord to increase the fanout once more. As Overlord does so, the average reliability increases too. As older messages are processed, reliability increases until it reaches 100\%. At minute 12, pit nodes are turned off. Coincidentally, three runs with 100\% reliability have passed and so Overlord decreases the fanout. This happens again, and so at minute 16, Overlord decreases the fanout once more. The rise in the reliability as a consequence of increasing the fanout is a positive result, although, these results do need to be compared to solutions not running Overlord in order to determine its benefits.

\begin{figure}
    \centering
	\includegraphics[scale=0.9]{FreeRiders/freeridersReliabilityComparisonWFanout.pdf}
	\caption{Comparison of Average Reliability Evolution over Time}
	\label{fig:freeridersReliabilityComparison1}
\end{figure}

Figure~\ref{fig:freeridersReliabilityComparison1} presents the evolution of the average reliability of the three different versions of the application. As it is visible, the depicted line corresponding to the application running Overlord, is the same as in Figure~\ref{fig:freeridersOverlordReliability1}.

\todo[inline]{I have a version of graph in Fig~\ref{fig:freeridersReliabilityComparison1} without the fanout decisions, is it better?}

All applications behave similarly during the first 4 minutes. When the pit nodes are enabled, both the application running Overlord and the application with the fanout set to 5, experience a dip in their average reliability, dropping to around 90\%. The optimistic application does not experience any drops in reliability during the system disturbance. The Overlord application, recovers due to increasing its fanout to account for the disturbance. In four metric collections, its average reliability returns to the 100\% mark, matching the pessimistic application.

\begin{figure}
    \centering
	\includegraphics[scale=0.9]{FreeRiders/freeridersLatencyComparison.pdf}
	\caption{Comparison of Average Latency Evolution over Time}
	\label{fig:freeridersLatencyComparison}
\end{figure}

Another aspect to consider, is the average latency of messages during the period in which pit nodes are enabled. As is presented in Figure~\ref{fig:freeridersLatencyComparison}, all three application behave similarly in the first 5 minutes with an average latency of around 200ms, and close to the 300ms for the Overlord and optimistic fanout applications. Upon the disturbance period, latency rises. However, for the optimistic application, it rises to over 600ms, while for the pessimistic application, only to around 350ms. In the case of Overlord, initially its latency rises to just under 600ms, while dropping with each increase of the fanout until it matches the 300ms latency reported by the pessimistic solution.

After these results, the question of why not simply use the pessimistic solution and have the fanout static at a higher value in case of disturbances arises. This version of the application, even though it is prepared for disturbances in the system, produces a lot of duplicate messages during phases of normal system functioning, which is also undesirable. Even though it suffers no dip in reliability, this comes at a cost, which is the high number of duplicate messages it produces. To measure this, we use the RMR metric, introduced in Section~\ref{par:rmr}.

Figure~\ref{fig:free_riders_rmr_comparison}, shows the comparison between the average RMR of each version of the application. Overlord adapts the fanout to achieve the best reliability value, with the lowest RMR\cite{Leitao2007a}. The positive performance of the application with a static fanout set to 8, comes at the cost of a high average RMR. The optimistic application, even though it presents a lower RMR, has the cost of a low reliability in times of system disturbance. The Overlord application, presents an alternative, in which during periods of normal functioning, the RMR is closer to the one presented by the optimistic application, and in times of system malfunction, it rises closer to the one presented by the pessimistic application, while providing an optimal reliability value in contrast to the optimistic application.

\begin{figure}
    \centering
	\includegraphics[scale=0.9]{FreeRiders/freeridersRMRComparison.pdf}
	\caption{Comparison of the Average RMR over Time.}
	\label{fig:free_riders_rmr_comparison}
\end{figure}

\subsubsection{Dealing With Loss of Information}
\label{sssec:tree_drop}
Our solution uses a disposable tree for collecting information. One aspect that surged during the elaboration of this dissertation was the uncertainty of the reaction from Overlord, when it loses information. As such, we decided to use the same scenario referred in Section~\ref{sssec:free_riders}, and compare different experiments with Overlord, where with each, more information would be lost. The objective is to find the percentage of the tree with which Overlord still reports a correct state of the system and still reacts to the disturbance in its environment.

\paragraph{Experiment Details} The details of this experiment are the same as in Section~\ref{sssec:free_riders}, however, instead of testing different static configurations, we only test Overlord. Moreover, during the period where pit nodes are enabled (which corresponds to 8 collections), in the first 3 collections, a certain percentage of the metric collection tree is lost. This provides us with a period where a disturbance is affecting system performance, and where Overlord does not have the complete information of the system. After these 3 collections the tree is restored. We tested four different versions of Overlord. The first one, did not lose any percentage of the tree, maintaining all information of the system at all times. Then, we tested versions where only 50\%, 25\% and 5\% of the nodes of the system remained. In the figures, yellow shaded regions represent collections where the information from nodes was lost, and the red shaded area represents the area where the system disturbance was still in place. The yellow area always overlaps the red one.

\begin{figure}
    \centering
	\includegraphics[scale=0.9]{TreeDrop/nodesEvolution.pdf}
	\caption{Evolution of the number of nodes over time.}
	\label{fig:tree_drop_nodes_evolution}
\end{figure}

Figure~\ref{fig:tree_drop_nodes_evolution} presents the evolution of the reported number of nodes overtime, for each of the tests. Note that these do not represent the amount of nodes present in the system, but only the number of nodes that successfully reported their information to the Overlord. A node being lost, means that it did not respond to its parent. This means that its child's information will also be compromised. Moreover, the exact number of nodes varies, as a new tree is created for each collection, meaning that different collections can have different trees. Nodes lost due to their parents in one collection, can have a different set of parents and be considered in the next one.

\begin{figure}
    \centering
	\includegraphics[scale=0.9]{TreeDrop/reliabilityEvolution.pdf}
	\caption{Evolution of average reliability over time.}
	\label{fig:tree_drop_reliability_evolution}
\end{figure}

\begin{figure}
    \centering
	\includegraphics[scale=0.9]{TreeDrop/fanoutEvolution.pdf}
	\caption{Evolution of average reliability over time.}
	\label{fig:tree_drop_fanout_evolution}
\end{figure}

Figure~\ref{fig:tree_drop_reliability_evolution} represents the evolution of the reported reliability by each of the applications. A way to know if the reported reliability is the correct one, is to observe if any dips occur when the information is restored (after the yellow shading). Since no new disturbances are introduced, there is no reason for a lower reliability, meaning simply that Overlord reported the wrong information.

Moreover, Figure~\ref{fig:tree_drop_fanout_evolution} represents the decisions made by Overlord. In addition to Figure~\ref{fig:tree_drop_reliability_evolution}, we can establish a reason for each decision of Overlord. First we discuss the version without any loss of information. As expected, it experiences a dip in reliability, which prompts the Overlord to increase the fanout, consequentially, improving it. The version where 50\% of the system is lost, is more complicated. After the pit nodes are enabled, this version reports a lower average reliability, which prompts the fanout to increase. However, instead of increasing to fanout 8, it stops at 7. This is due to the cap we introduced to the fanout increments. As Overlord believes there are only around 100 nodes (instead of the 200), it can only increase its fanout to this value. However, this has a positive effect, as after information is restored, reliability is above the accepted threshold. The version with 25\% of the system also increases its fanout to 7, however after the information is restored, the reliability reported is sightly lower than the previous one. This indicates us that Overlord did not have a correct value of the real reliability, however, it quickly recovers, increasing its fanout to 8 and restoring the reliability. Finally, the version with only 5\% of the system, has a completely incorrect assessment of the reliability, as it reports 80\%, which rises to almost 90\%. After information is restored, the real reliability is presented, which corresponds to 40\%.

In essence, losing information, always affects the decisions Overlord can make, especially due to the cap on reliability that we introduced. As such, decisions should not be affected, as long as the amount of nodes remaining does not affect the optimal fanout. In terms of reported information, the application with 50\% of the nodes correctly depicted the true reliability of the system, whereas the remaining versions showed gaps between the information monitored and the true state of the system, with the 5\% version of the system, having a significantly wrong idea of the system's performance. 

\subsubsection{Packet Loss}
\label{sssec:packet_loss}

This experiment consists of deliberately dropping a certain percentage of the messages broadcasted by the nodes in the system, meaning that for each message a node receives, there is a percentage $p$ that this message will be lost and not delivered. This will lower the reliability because some messages will not be registered by all nodes, which should prompt Overlord to increase the fanout in order to improve the reliability. Increasing the fanout should work since it means increasing the number of neighbors to which a node sends a message, which increases the number of messages sent.

Figure~\ref{fig:packetlossexperiment3} shows the main principle behind this experiment, by representing a small portion of a hypothetical system. In Figure~\ref{fig:packetlossexperiment1}, node $a$ creates message $m$. Since fanout is set to 4, the message is sent to all four neighbors ($b$, $c$, $d$ and $e$), but since 4 out of 5 messages are lost, none of them actually deliver the message. If the fanout increases to 5, as is represented in Figure~\ref{fig:packetlossexperiment2}, $f$ is added to $a$'s membership, meaning that $f$ will receive message $m$. Node $f$ will also send the message to 5 of its neighbors, but it will only be delivered by one of them. If the fanout increases again, 2 neighbors receive the message, and this can happen until enough nodes receive the message in order to disseminate it to the whole system.

Since Babel is built on TCP connections, lost messages is not truly a realistic scenario. If a message is lost, retransmission is triggered until the message is delivered. In this application, a message is only registered as delivered if it is delivered during a certain period of time, so even if a node receives a retransmitted message, it may not be delivered in time\cite{Lee2002}. As such, this experiment is not truly a realistic scenario, and was created as a methodical and more deterministic way of lowering the reliability of the system on demand, and observe the effects of adapting the fanout parameter\todo{Here choose if we can justify this experiment being real or if it is just a "lower reliability" experiment}. Instead of devising a more realistic different scenario, which would make observations between comparisons of different solutions more difficult\todo{What more experiments could we have done?}, this experiment is a simple disturbance that reduces the number of biases that can happen during testing, and helps observe the differences between having a static fanout, and having a fanout adapted to disturbances.

\begin{figure}
\centering
\begin{subfigure}{.46\textwidth}
  \centering
  \includegraphics[scale=2.1]{PacketLoss/packetLossA.pdf}
  \caption{Since 4 out of 5 messages are lost, none of the neighbors of $a$ receive message $m$.}
  \label{fig:packetlossexperiment1}
\end{subfigure}%
\hspace{1cm}
\begin{subfigure}{.46\textwidth}
  \centering
  \includegraphics[scale=2.1]{PacketLoss/packetLossB.pdf}
  \caption{As the fanout increased to 5, $a$ now has a new neighbor $f$ receives $m$.}
  \label{fig:packetlossexperiment2}
\end{subfigure}
\caption{A visual example of the expected improvements of Overlord on the packet loss experiment.}
\label{fig:packetlossexperiment3}
\end{figure}

\paragraph{Experiment Details} \label{par:packetlossexperiment_details} The experiment was conducted with 200 nodes participating in the system, with the threshold of acceptable reliability at 98\%. As in the other experiments, we will be comparing three versions of the application \begin{inparaenum}[(i)] \item A version with an optimistic static fanout, which does not consider any disturbances on the system, \item A version with a pessimistic static fanout, considering the consequences of this experiment, and \item A version running Overlord and the rule specified in Section~\ref{ssec:fanout_adaptation_process}, which will adapt the fanout in accordance to the monitored reliability\end{inparaenum}. During the first 5 metric collections, the system works as intended, without any disturbances in order to simulate normal operation. At the 5th metric collection, packet dropping is turned on, and every message received from other nodes has a 25\% chance of being discarded. This behavior lasts for 8 more rounds, in order to allow Overlord to react to any issues found, and at the end of these 8 rounds, normal behavior is restored.

Similarly to the free riders experiment, HyParView can introduce a certain bias, making it harder to compare different versions of the same application due to the membership changing. In the same way as mentioned in Section~\ref{par:freeridersexperiment_details}, the same set of HyParView generated memberships was used for every version of the application tested in order to reach a less biased conclusion.

\begin{figure}
    \centering
	\includegraphics[scale=0.9]{PacketLoss/packetlossOverlordReliability.pdf}
	\caption{Average reliability evolution over Time}
	\label{fig:packetlossOverlordReliability}
\end{figure}

Figure~\ref{fig:packetlossOverlordReliability} presents the fluctuations of the average reliability over time monitored by Overlord, as well as the different decisions it made on each metric collection. The red shaded area represents the period in which packets were being lost. During the first 4 minutes, the system runs as expected with no disturbance, during 5 metric collections. At the 5th collection, the packet loss is enabled in all nodes, and packets are lost with a probability of 25\%. As observable in Figure~\ref{fig:packetlossOverlordReliability}, the first dip occurs at the 5-minute mark with an average reliability of 87\%. Due to some messages that were sent before packet loss activation still being processed, the reliability is not yet at it lowest. However, it is low enough for Overlord to increase the fanout. At minute 6, the average reliability is at its lowest, with a value of 76\%, meaning that Overlord increases the fanout to 7. This results in an improvement to the reliability, leading it to rise to 87\%, which is still below the acceptable threshold. Consequently, Overlord increases the fanout to 8, which prompts the reliability to increase to around 92\%. This value is still under the acceptable threshold of 98\%, however, Overlord cannot increase it again due to having reached the maximum value for the fanout parameter defined for this experiment, which is $\log{200} + 3 = 8$. This cap prevents the fanout growing infinitely. If this value were adjusted, Overlord would increase it again, as it was depicted in the free riders experiment, Section~\ref{sssec:free_riders}.

\todo[inline]{Maybe this isn't good. Because I'm essentially saying that Overlord does not increase the fanout on its own to get the best reliability. I have tests with a lower packet drop percentage (15\%) which don't present the same problems. We have to decide if it's better to show a "feature" of the rule in order for it not to grow infinitely, or if It's too much of a "My solution doesn't work as well as it should because I limited it", and use a more perfect version of the test.}

As mentioned previously, these results alone are not enough to draw conclusions, so two other versions of the application were also tested in order to compare these different options, and determine what differences using Overlord provides. These different versions, are the ones referred in Section~\ref{ssec:evaluation_scenarios}.

\begin{figure}
    \centering
	\includegraphics[scale=0.9]{PacketLoss/packetlossReliabilityComparisonWFanout.pdf}
	\caption{Comparison of Average Reliability Evolution over Time}
	\label{fig:packetlossReliabilityComparison}
\end{figure}

\todo[inline]{I have a version of graph in Fig~\ref{fig:packetlossReliabilityComparison} without the fanout decisions, is it better?}

Figure~\ref{fig:packetlossReliabilityComparison} shows the evolution of the average reliability over time of the three different versions of the application that were tested in this experiment. The overlord curve is the same as shown in Figure~\ref{fig:packetlossOverlordReliability}. During the normal functioning period of the application, all three applications register a similar average reliability of 100\%. The first collection after enabling the packet loss shows that both the Overlord and the optimistic application experience a similar dip in their average reliability, falling to around 87\%. At the 7th minute, reliability keeps lowering, however, the optimistic application experiences a lower reliability of 70\%, in contrast to the Overlord application, which reports a reliability of 76\%. This is due to Overlord having changed its fanout to 6. This distancing from the optimistic application continues as Overlord increases its fanout (as shown in Figure~\ref{fig:packetlossOverlordReliability}), until it matches the pessimistic application with an average reliability of around 92\%. The optimistic application remains in the 70\% reliability for the entirety of the packet loss period. Consequentially, for a single message, out of the 200 nodes participating in the system, about 60 nodes do not receive that message.

\begin{figure}
    \centering
	\includegraphics[scale=0.9]{PacketLoss/packetlossLatencyComparison.pdf}
	\caption{Comparison of Average Latency Evolution over Time}
	\label{fig:packetlossLatencyComparison}
\end{figure}

Moreover, it is also important to note that the latency during the disturbance period increases significantly for the optimistic fanout application, going from around 200ms in periods of normal functioning, to over 700ms during packet loss, as shown in Figure~\ref{fig:packetlossLatencyComparison}. The pessimistic application also experiences a higher latency, however much lower than the optimistic application, with a latency of 400ms. Overlord, initially follows the optimistic application in terms of latency values, but as it increases its fanout, this value drops from 600ms in the second metric collection during the packet loss period, to match the pessimistic application around 400ms.

\begin{figure}
    \centering
	\includegraphics[scale=0.9]{PacketLoss/packetlossRMRComparison.pdf}
	\caption{Comparison of Average Reliability Evolution over Time}
	\label{fig:packet_loss_rmr_comparison}
\end{figure}

Similarly to the free riders experiment, simply using a static pessimistic fanout is not a viable option as it produces a higher RMR than needed during the normal functioning period of the test. As presented by Figure~\ref{fig:packet_loss_rmr_comparison}, during the first 5 minutes of the experiment, the pessimistic application registers an average RMR of around 6, while presenting the same average reliability as the other two applications. Moreover, the pessimistic application has a static fanout relatively to the 25\% of packets being lost. If this value were higher, this version of the application would have a worse performance, whereas Overlord would adapt the fanout to achieve a better reliability value.

\section{Summary}
\label{sec:evaluation_summary}
In this chapter, we evaluated the developed architecture, implementing a message dissemination application as a case study on top of the created framework. First we explain the application, giving an overview from the "point of view" of a developer using Overlord to monitor and adapt this system. Then, we presented the impact of the average reliability of the system, which helps understand system performance, since it directly indicates the percentage of nodes that received messages. We also detailed how Overlord uses the rules tailored to this system in order to adapt the fanout parameter according to the average reliability monitored in each metric collection. Finally, we evaluated how Overlord performed in adapting this application by testing it in different possible scenarios that can occur during execution of these systems. Evaluation was performed by testing the adaptive solution against optimistic and pessimistic applications, with a static fanout. Our solution showed that it is capable of adapting to periods of instability, not succumbing to such issues as the optimistic application, while also being capable of optimizing itself during periods of positive performance, not introducing more duplicate messages than needed as does the pessimistic application.